\newcommand{\figuraFuerzasMolecularesSuperficie}{%
\begin{figure}[H]
\centering
\begin{tikzpicture}[>=Latex,font=\small,line cap=round,line join=round]
  \fill[blue!10] (0.7,0.5) rectangle (11.3,4.2);
  \draw[very thick,blue!65!black] (0.7,4.2)--(11.3,4.2);
  \foreach \x in {1.4,2.4,...,10.4}
    \foreach \y in {1.1,2.1,3.1,4.1}
      \fill[blue!60] (\x,\y) circle (0.13);
  \coordinate (I) at (3.4,2.1);
  \coordinate (S) at (8.4,4.1);
  \foreach \a in {0,45,...,315}
    \draw[->,thin,green!50!black] (I)--++(\a:0.65);
  \foreach \a in {180,225,270,315,360}
    \draw[->,thin,red!75!black] (S)--++(\a:0.68);
  \draw[->,ultra thick,purple!75!black] (S)--++(0,-1.25)
    node[midway,right] {resultante};
  \node[align=center] at (3.4,5.0)
    {molécula interior:\\[1mm]fuerzas compensadas};
  \node[align=center] at (8.4,5.0)
    {molécula superficial:\\[1mm]atracción hacia el líquido};
\end{tikzpicture}
\caption{Origen molecular de la tensión superficial. En el interior las
atracciones se compensan; junto a la interfaz faltan vecinas del lado gaseoso
y aparece una tendencia neta a reducir el área.}
\label{fig:fuerzas-moleculares-superficie}
\end{figure}%
}